\begin{abstract}
Although large unsupervised language models (LMs) acquire extensive world knowledge and some capacity for reasoning, their behavior remains challenging to regulate precisely because their training is entirely unsupervised. To improve controllability, current approaches rely on human-annotated comparisons between model outputs, subsequently fine-tuning the LM to reflect these labeled preferences—commonly utilizing reinforcement learning from human feedback (RLHF). Nevertheless, RLHF involves a multi-stage and frequently brittle process: it first constructs a reward model embodying human preferences, then adapts the original LM through reinforcement learning to maximize the estimated reward, all while striving to prevent significant divergence from the pre-trained model. In this work, we propose a novel formulation of the reward model within RLHF that makes it possible to directly recover the optimal policy in closed form, thereby reducing the standard RLHF objective to a simple classification loss. Our proposed method, termed \textit{Direct Preference Optimization} (DPO), offers a stable, efficient, and lightweight alternative—eschewing the need for language model sampling during fine-tuning and minimizing extensive hyperparameter adjustment. Experimental results demonstrate that DPO is at least as effective as prior approaches for aligning LMs with human preferences, if not superior. Importantly, DPO outperforms PPO-based RLHF in guiding the sentiment of generated texts and matches or surpasses existing techniques for tasks like summarization and single-turn dialogue, all while streamlining both the implementation and training process.
\end{abstract}

\section{Introduction}
Large-scale unsupervised language models (LMs) trained on massive corpora exhibit an array of unexpected capabilities~\citep{chowdhery2022palm, brown2020language, touvron2023llama, bubeck2023sparks}. Nonetheless, these models are built on data sourced from humans with a diverse range of intentions, priorities, and expertise. Not all of these human tendencies are necessarily beneficial to replicate; for instance, it is desirable for an AI programming assistant to \textit{recognize} frequent coding errors so it can fix them, but during code generation, we prefer the system to favor the (sometimes uncommon) high-level programming proficiency represented in its training material. In a similar vein, a language model should be \textit{informed} about widespread misconceptions—for example, one held by half of the population—but we would not want it to reproduce this misconception as a correct answer in half its responses. Thus, carefully \emph{aligning model outputs and behaviors} with intended objectives, as distinct from their broad \textit{capabilities and information}, is essential for developing safe, effective, and controllable AI~\citep{ouyang2022training}. Although standard practice aligns LMs to human preferences via reinforcement learning (RL), we demonstrate that the RL-based loss function commonly used in such approaches can be solved exactly through a straightforward binary cross-entropy objective, thus streamlining the process of preference alignment.

Broadly, current alignment strategies involve endowing language models with desired behavioral traits by leveraging meticulously constructed datasets that reflect human notions of safety and helpfulness. This preference adaptation phase follows extensive unsupervised pre-training on large textual datasets. While the simplest preference alignment strategy is to fine-tune using human-annotated examples of high-quality responses, the most effective approach to date is reinforcement learning from human (or artificial) feedback (RLHF/RLAIF; \citep{christiano2017deep, bai2022constitutional}). RLHF involves training a reward model based on preference data and subsequently using RL techniques to tune a language model to generate responses deemed high-reward, while constraining deviations from the initial model distribution. Although RLHF has led to highly capable models for tasks such as conversation and code generation, it introduces substantial complexity compared to supervised approaches: the procedure requires the training of several language models and the repeated sampling from the model’s policy within the training loop, resulting in significantly increased computational requirements.

This work introduces a method for optimizing language models according to human preferences without relying on explicit reward modeling or reinforcement learning techniques. We present \textit{Direct Preference Optimization (DPO)}, a novel approach that implicitly targets the same objective as established RLHF methods—namely, reward maximization under a KL-divergence regularization—but with an implementation and training process that is notably more accessible. In essence, DPO updates operate by enhancing the log probability ratio between preferred and non-preferred outputs; crucially, it integrates an adaptive, per-example importance weighting scheme to avoid the model collapse seen with simplistic probability ratio formulations. Analogous to prior approaches, DPO utilizes a theoretical preference model (e.g., the Bradley-Terry model; \cite{bradley1952rankanalysis}) to quantify the alignment between a reward function and observed preference data. Unlike these methods, which train a reward model via preference loss and then optimize a policy against this learned reward, DPO reparameterizes the preference loss to depend directly on the policy itself. Utilizing human preference annotations over model outputs, DPO enables straightforward policy training with a binary cross entropy objective, leading to the optimal policy for an implicit reward inferred from the preference information.

The principal contribution of this study is Direct Preference Optimization (DPO), a streamlined, reinforcement learning-free technique for preference-driven language model training. Empirically, we demonstrate that DPO performs on par with leading alternatives, including PPO-based RLHF, across a range of tasks such as sentiment adjustment, summarization, and conversational modeling, with language models scaling up to 6B parameters.

Large self-supervised language models increasingly demonstrate the capability to solve various tasks in a zero-shot manner \citep{radford2019language} or through the use of few-shot prompting techniques \citep{gpt3,megatron,chowdhery2022palm}. Nonetheless, their downstream task accuracy and alignment with human intentions can be further enhanced through fine-tuning on collections of task instructions and human-authored completions \citep{mishra-etal-2022-cross,sanh2022multitask,chung2022scaling,thoppilan2022lamda}. This approach, known as `instruction-tuning,' allows large language models (LLMs) to generalize more effectively to novel instructions not present in the fine-tuning dataset, thus improving overall model utility \citep{chung2022scaling}. Although instruction-tuning has proven effective, obtaining human rankings of response quality (i.e., relative preferences) is often less demanding than gathering expert-crafted demonstrations. Consequently, recent studies have focused on further fine-tuning LLMs with preference-based datasets, resulting in improved outcomes for tasks such as translation \citep{kreutzer-etal-2018-reliability}, summarization \citep{stiennon2022learning,ziegler2020finetuning}, narrative generation \citep{ziegler2020finetuning}, and adherence to instructions \citep{ouyang2022training,ramamurthy2023is}. The common methodology involves first training a neural reward model to align with human preference data, typically modeled using approaches like the Bradley-Terry model \citep{bradley1952rankanalysis}. Subsequently, the language model is fine-tuned to optimize this learned reward function via reinforcement learning algorithms such as REINFORCE \citep{williams1992reinforce}, proximal policy optimization (PPO; \cite{schulman2017proximal}), or their adaptations \citep{ramamurthy2023is}. Related research uses instruction-following LLMs trained with human feedback to generate additional synthetic preference labels emphasizing particular characteristics, such as safety or harmlessness \citep{bai2022constitutional}, relying only on minimal human oversight in the form of rubric-like guidance for annotation. These methodologies sit at the intersection of two research streams: one focusing on reinforcement learning approaches for language models across multiple objectives~\citep{Ranzato2015SequenceLT,paulus2018a,wu2018learning}, and another dedicated to the general problem of learning from human preferences \citep{christiano2017deep,kupcsik2018learning}. Despite the promise of leveraging relative preferences for model fine-tuning, employing reinforcement learning to adjust large language models presents considerable practical difficulties; this work introduces a theoretically-motivated technique for optimizing based on relative preferences that circumvents the use of reinforcement learning.

Beyond language-focused scenarios, the study of policy learning from preferences has a history in both bandit frameworks and reinforcement learning, with multiple methodologies proposed. In contextual bandit settings, leveraging preferences or action rankings instead of explicit reward signals is formalized as the contextual dueling bandit (CDB; \cite{yue2012karmed,dudik2015contextual}) problem. When absolute rewards are unavailable, analysis of CDBs replaces the standard concept of policy optimality with the \textit{von Neumann winner}, defined as a policy whose probability of outperforming \textit{every} other policy is no less than 50\% on average \citep{dudik2015contextual}. A key distinction of the CDB formulation is that preferences are supplied in an online fashion, whereas human preference learning often operates over a static, offline collection of action pairs annotated with preference information \citep{yan2022human}. In the same vein, \textit{preference-based RL} (PbRL) trains on binary preference labels drawn from an \textit{unknown} scoring mechanism, as opposed to observable rewards \citep{BusaFekete2014,ruiz2023dueling}. The PbRL literature offers several algorithms, some of which are designed to leverage off-policy preference data. Typically, these methods follow a two-step pipeline: first, the underlying (and latent) scoring function is estimated—commonly understood as learning the reward model—followed by its subsequent optimization \citep{jain2013learning,BusaFekete2014,christiano2017deep,sadigh2017active,kupcsik2018learning}. Our work instead departs from this approach by introducing a single-phase policy learning procedure that directly tailors a policy to adhere to provided preference signals.

\section{Preliminaries}\label{section:prelims}

We summarize the RLHF methodology as outlined by \citeauthor{ziegler2020finetuning}, subsequently expanded in \citep{stiennon2022learning, bai2022training, ouyang2022training}. The standard pipeline encompasses three distinct steps: 1) supervised fine-tuning (SFT); 2) collecting preference data and fitting a reward model; and 3) optimization using reinforcement learning.

\textbf{SFT}: The process commences with adapting a pre-trained language model to the target domain(s)—such as dialogue or summarization—via supervised training on curated, high-quality datasets, resulting in the SFT model $\pi^\text{SFT}$.

\textbf{Reward Modelling Phase}: During this step, the fine-tuned model receives prompts $x$, generating pairs of responses $(y_1, y_2)\sim \pi^\text{SFT}(y \mid x)$. Human annotators are then presented with these completions and express a preference, written $y_w\succ y_l \mid x$, indicating $y_w$ is favored over $y_l$ for a given $x$. The source of these preferences is presumed to be an unknown reward function $r^*(y, x)$ to which we do not have explicit access. Several techniques have been developed to capture and model these preference signals, with the Bradley-Terry (BT) \cite{bradley1952rankanalysis} framework being prevalent. More broadly, ranking models such as Plackett-Luce \citep{plackett1975analysis, luce2012individual} can also be incorporated if sets of ranked responses are available. Under the BT formalism, the distribution $p^*$ of human preferences is specified as follows:

\begin{equation}\label{eq:bradley-terry}
    p^*(y_1\succ y_2 \mid x)=\frac{\exp\left(r^*(x, y_1)\right)}{\exp\left(r^*(x, y_1)\right) + \exp\left(r^*(x, y_2)\right)}.
\end{equation}

Suppose we are given a fixed dataset of comparative examples, denoted as $\mathcal{D}=\bigl\{x^{(i)}, y_w^{(i)}, y_l^{(i)}\bigr\}_{i=1}^N$, which are drawn according to $p^*$. We introduce a parameterized reward model $r_{\phi}(x, y)$ and fit its parameters by maximizing the likelihood of the observed data. When cast as a binary classification task, this leads to minimizing the negative log-likelihood loss defined as follows:

\begin{equation}\label{eq:reward_model}
    \mathcal{L}_R(r_{\phi}, \mathcal{D}) = -\mathbb{E}_{(x, y_w, y_l)\sim \mathcal{D}}\bigl[\log \sigma(r_{\phi}(x, y_w)- r_{\phi}(x, y_l))\bigr]
\end{equation}

with $\sigma$ denoting the logistic sigmoid function. In language modeling contexts, $r_{\phi}(x, y)$ is frequently initialized from the SFT model $\pi^\text{SFT}(y \mid x)$, appending a linear projection atop the transformer’s last hidden layer to yield a scalar reward estimate \cite{ziegler2020finetuning}. To mitigate reward function variance, it is standard practice to normalize the output so that $\mathbb{E}_{x,y\sim \mathcal{D}}\left[r_\phi(x, y)\right] = 0$ for every $x$.

\textbf{RL Fine-Tuning Phase}: During the reinforcement learning stage, the model uses the trained reward signal to update the language model. Consistent with earlier research~\citep{jaques2017sequence, jaques2020human}, the optimization procedure is formulated as

\begin{equation}\label{eq:RL}
\max_{\pi_{\theta}}  \mathbb{E}_{x\sim \mathcal{D}, y\sim \pi_{\theta}(y \mid x)}\bigl[r_{\phi}(x, y)\bigr] - \beta\mathbb{D}_{\textrm{KL}}\bigl[\pi_{\theta}(y\mid x)\mid \mid \pi_\text{ref}(y\mid x)\bigr],
\end{equation}

where the parameter $\beta$ governs how closely $\pi_{\theta}$ adheres to a reference policy, $\pi_\text{ref}$, usually set to the SFT initialization $\pi^\text{SFT}$. Typically, $\pi_\theta$ is also initialized with the parameters of $\pi^\text{SFT}$. The regularization term constrains policy updates to stay near the distribution where the reward model remains reliable, promotes output variety, and reduces the risk of collapsing to only high-reward responses. Due to the inherent discreteness of sequence generation, the optimization objective is non-differentiable and must generally be addressed via reinforcement learning techniques. The predominant method \citep{ziegler2020finetuning, stiennon2022learning, bai2022training, ouyang2022training} is to use a reward shaped as ${r(x, y) = r_{\phi}(x, y) -\beta (\log \pi_{\theta}(y\mid x) - \log \pi_\text{ref}(y\mid x))}$ and apply algorithms such as PPO \cite{schulman2017proximal} for maximization.

\section{Direct Preference Optimization}\label{sec:DPO}

Recognizing the complexity inherent in deploying RL algorithms at the scale required for large language model tuning, we seek to develop a streamlined approach to policy optimization that utilizes preference data directly. Unlike established RLHF methodologies, which decouple the reward learning from policy optimization, our proposed technique employs a special choice of reward parameterization which makes it possible to compute the optimal policy analytically, thereby circumventing the need for reinforcement learning loops. As detailed below, our main contribution is to utilize a closed-form transformation between reward functions and optimal policies, enabling us to convert an objective over reward models into an objective defined directly on policies. This variable transformation dispenses with the necessity for an explicit reward model, while still adhering to established preference models such as the Bradley-Terry framework. Consequently

\textbf{Derivation of the DPO objective.} We begin with the standard reinforcement learning objective given in Eq.~\ref{eq:RL}, involving an arbitrary reward function $r$. Consistent with previous works~\citep{peters2007reinforcement, peng2019advantage, korbak2022reinforcement, go2023aligning}, it is straightforward to demonstrate that the optimal policy for the KL-regularized reward maximization problem in Eq.~\ref{eq:RL} adopts the following form:

\begin{equation}\label{eq:op_policy}
    \pi_r(y\mid x) = \frac{1}{Z(x)}\pi_\text{ref}(y\mid x)\exp\left(\frac{1}{\beta}r(x, y)\right),
\end{equation}

with $Z(x) =\sum_{y}\pi_\text{ref}(y\mid x)\exp\left(\frac{1}{\beta}r(x, y)\right)$ serving as the normalization constant, or partition function. Details on this derivation are provided in Appendix \ref{app:derivation1}. Even when utilizing an MLE-trained reward estimator $r_{\phi}$ to approximate the true reward $r^*$, calculating $Z(x)$ remains computationally intractable in practice~\citep{korbak2022reinforcement, go2023aligning}. This limitation complicates the direct application of this optimal policy representation. However, by rearranging Eq.~\ref{eq:op_policy}, the reward can be expressed as a function of its optimal policy $\pi_r$, the reference $\pi_\text{ref}$, and the unknown $Z(\cdot)$. Taking the logarithm of Eq.~\ref{eq:op_policy} and reordering terms gives:

\begin{equation}\label{eq:main_eq}
    r(x,y) =\beta \log \frac{\pi_r(y\mid x)}{\pi_\text{ref}(y\mid x)} + \beta \log Z(x).
\end{equation}

This form allows substitution for the true reward $r^*$ and the optimal policy $\pi^*$. Notably, in the Bradley-Terry preference framework, only the reward difference between two outputs is required: ${p^*(y_1 \succ y_2 \mid x) = \sigma(r^*(x, y_1) - r^*(x, y_2))}$. Inserting the parameterization from Eq.~\ref{eq:main_eq} for $r^*(x,y)$ into the preference expression (Eq.~\ref{eq:bradley-terry}), the partition function terms negate each other, enabling us to represent human preference probabilities purely in terms of the optimal policy $\pi^*$ and reference policy $\pi_\text{ref}$. Therefore, the policy $\pi^*$ that is optimal for RLHF under the Bradley-Terry assumption satisfies:

\begin{equation}\label{eq:objective}
    p^*(y_1\succ y_2 \mid x)=\frac{1}{1 + \exp\left(\beta \log \frac{\pi^*(y_2\mid x)}{\pi_\text{ref}(y_2\mid x)} - \beta \log \frac{\pi^*(y_1\mid x)}{\pi_\text{ref}(y_1\mid x)}\right)}
\end{equation}

The steps leading to this result are presented in Appendix~\ref{app:derivation2}. While Eq.~\ref{eq:objective} specifically employs the Bradley-Terry model, similar derivations apply for the more general Plackett-Luce family~\citep{plackett1975analysis, luce2012individual}, as detailed in Appendix~\ref{app:plackett_luce_models}.

With this representation, the probability of human preferences is written directly in terms of the policy, avoiding explicit reward computation. This enables us to define a maximum likelihood criterion for a policy parameterization $\pi_\theta$, analogous to reward-modeling methods (see Eq.~\ref{eq:reward_model}). Specifically, the resulting objective for DPO is:

\begin{equation}\label{eq:optimum_model}
    \mathcal{L}_\text{DPO}(\pi_{\theta}; \pi_\text{ref}) = -\mathbb{E}_{(x, y_w, y_l)\sim \mathcal{D}}\left[\log \sigma \left(\beta \log \frac{\pi_{\theta}(y_w\mid x)}{\pi_\text{ref}(y_w\mid x)} - \beta \log

In this approach, we model the implicit reward function via an alternative parameterization, with the resulting optimal policy being $\pi_\theta$. Additionally, as our method aligns with a reparameterized Bradley-Terry framework, it inherits several desirable theoretical guarantees—such as consistency, given appropriate conditions on the distribution of preference data \cite{bong2022generalized}. Further theoretical analysis of DPO, including its connections to prior literature, is provided in Section~\ref{sec:theory}.

\textbf{How does the DPO update work?} To gain a mechanistic perspective on DPO, one can examine the gradient of the objective $\mathcal{L}_\text{DPO}$. The parameter gradient is given by:

\begin{multline*}\label{eq:gradient}
    \nabla_\theta \mathcal{L}_\text{DPO}(\pi_\theta;\pi_\text{ref}) = \\ -\beta\mathbb{E}_{(x, y_w, y_l) \sim \mathcal{D}} \bigg[\underbrace{\sigma(\hat{r}_\theta(x, y_l) - \hat{r}_\theta (x, y_w))}_\text{increased emphasis when reward ordering is inaccurate}\bigg[\underbrace{\nabla_\theta\log \pi(y_w \mid x)}_\text{boost $y_w$ probability} - \underbrace{\nabla_\theta\log\pi(y_l \mid x)}_\text{lower $y_l$ probability}\bigg]\bigg],
\end{multline*}

where the implicit reward function is defined as $\hat{r}_\theta(x, y) = \beta \log \frac{\pi_\theta(y \mid x)}{\pi_\text{ref}(y \mid x)}$, constructed from both the language model $\pi_\theta$ and the reference model $\pi_\text{ref}$ (with further discussion in Section~\ref{sec:theory}). In essence, the loss gradient $\mathcal{L}_\text{DPO}$ promotes higher probabilities for preferred responses $y_w$ while diminishing probabilities for less-preferred $y_l$ outcomes. Crucially, the influence of each training example is scaled according to how much the current reward model $\hat{r}_\theta$ overestimates dispreferred completions—reflecting the extent of its ranking error—modulated by $\beta$ to incorporate the strength of the KL regularization. Our empirical findings underscore the importance of this scaling, as omitting it leads to model collapse (see Appendix Table~\ref{tab:unlikelihood_generations} for details).

\textbf{DPO Overview.} The typical DPO workflow consists of the following steps: 1) For each prompt $x$, draw two completions $y_1, y_2$ from $\pi_\text{ref}(\cdot \mid x)$, annotate them according to human preferences, and assemble the resulting data into an offline preference set $\mathcal{D} = \{x^{(i)}, y_w^{(i)}, y_l^{(i)}\}_{i=1}^N$; 2) Train the language model $\pi_\theta$ by minimizing $\mathcal{L}_\text{DPO}$, conditioned on the specified $\pi_\text{ref}$, dataset $\mathcal{D}$, and selected value of $\beta$. In practice, practitioners often wish to take advantage of publicly available preference datasets rather than repeatedly generating completions and collecting fresh human judgments. When these preference datasets have been curated using $\pi^\text{SFT}$, we set the reference model $\pi_\text{ref} = \pi^\text{SFT}$ where possible. In instances where $\pi^\text{SFT}$ cannot be accessed, we instead fit $\pi_\text{ref}$ by maximizing the likelihood on preferred pairs $(x, y_w)$: specifically, ${\pi_\text{ref} = \argmax_{\pi}\mathbb{E}_{x, y_w \sim \mathcal{D}}\left[\log \pi(y_w \mid x)\right]}$. This method serves to reduce discrepancies between the unavailable true reference distribution and the operational $\pi_\text{ref}$ utilized during DPO. For specifics regarding the implementation and the chosen hyperparameters, please consult Appendix~\ref{app:implementation}.

\section{Theoretical Analysis of DPO} This section provides deeper insight into the DPO algorithm by offering a theoretical interpretation, examining its foundational guarantees, and highlighting the benefits of DPO relative to actor-critic algorithms commonly employed in RLHF (e.g., PPO~\cite{schulman2017proximal}).

\label{sec:theory}

\subsection{Your Language Model Is Secretly a Reward Model} DPO achieves the policy optimization goal by transforming the problem into a single maximum likelihood objective, eliminating the need for an explicit reward function or separate RL procedures. The loss function presented in Eq. \ref{eq:main_eq} corresponds mathematically to a Bradley-Terry model, wherein the reward is parameterized by $r^*(x, y) = \beta \log\frac{\pi^*_\theta(y \mid x)}{\pi_\text{ref}(y \mid x)}$. Thus, optimizing the model $\pi_{\theta}$ can be reinterpreted as optimizing a reward model, as in Eq. \ref{eq:reward_model}, once a variable transformation is made. In what follows, we will develop the theoretical framework for this reparameterization, demonstrate that it preserves the generality of the set of potential reward models, and show that it enables precise recovery of the optimal policy. We begin by introducing an equivalence relation for reward functions.

\begin{definition}
Two reward functions $r(x, y)$ and $r'(x, y)$ are considered equivalent if and only if there exists a function $f$ such that ${r(x, y)-r'(x, y) = f(x)}$.
\end{definition}

This equivalence relation readily partitions the space of reward functions into distinct equivalence classes. With this framework, we can establish the following two lemmas:

\begin{lemma}\label{lemma:same_prefrence}
Within the Plackett-Luce model, and specifically under the Bradley-Terry formulation, any pair of reward functions belonging to the same equivalence class yields identical preference distributions.
\end{lemma}

\begin{lemma}\label{lemma:same_policy}
    If two reward functions reside within the same equivalence class, they induce the same optimal policy in the corresponding constrained RL problem.
\end{lemma}

We note that the proofs of these statements are straightforward and can be found in Appendix~\ref{app:lemma1}. The first lemma echoes the classical under-determination encountered in the Plackett-Luce family of preference models \cite{plackett1975analysis}. As a consequence of this ambiguity, identifiability constraints are commonly enforced to enable meaningful estimation of the reward parameters via maximum likelihood in Eq.~\ref{eq:reward_model} \cite{bong2022generalized}. According to the second lemma, every reward function within an equivalence class produces the same optimal policy, indicating that for practical purposes it suffices to recover any representative from the correct class. The following theorem, which we demonstrate in Appendix~\ref{app:thm1}, captures this insight:

\begin{theorem}\label{thm:main}
    Assuming mild regularity conditions, each class of reward functions compatible with Plackett-Luce (and in particular, Bradley-Terry) models can be expressed as ${r(x, y) = \beta \log \frac{\pi(y\mid x)}{\pi_\text{ref}(y\mid x)}}$ for some model $\pi(y\mid x)$ and a specified reference model $\pi_\text{ref}(y \mid x)$.
\end{theorem}

\begin{sproof}
    Let $r(x, y)$ denote a reward function that determines an associated optimal policy $\pi_r(y \mid x)$ as given by Eq.~\ref{eq:op_policy}. To demonstrate that every function in the equivalence class of $r$ admits the above reparameterization, define the transformation $f$ by  
\begin{equation}
    f(r; \pi_\text{ref}, \beta)(x, y) = r(x, y) - \beta\log\sum_{y}\pi_\text{ref}(y\mid x)\exp\left(\frac{1}{\beta}r(x, y)\right)
\end{equation}
This transformation $f$ normalizes the reward function by subtracting the log-partition function of $\pi_r$, which depends solely on $x$. Therefore, $f(r; \pi_\text{ref}, \beta)(x, y)$ remains in the equivalence class of $r(x, y)$. Furthermore, by substituting $r$ as specified by Eq.~\ref{eq:main_eq} (which is valid for all reward functions), we find that $f(r; \pi_\text{ref}, \beta)(x, y) = \beta \log \frac{\pi_r(y\mid x)}{\pi_\text{ref}(y\mid x)}$. Consequently, the mapping $f$ yields a member of the equivalence class of $r$ possessing the required form, showing that our reparameterization retains the full generality necessary for the reward model.
\end{sproof}

Theorem~\ref{thm:main} can also be interpreted as identifying the specific reward function selected by the DPO reparameterization within each equivalence class; namely, the reward function that fulfills:

\begin{equation}\label{eq:lag_p}
     \sum_{y}\underbrace{\pi_\text{ref}(y\mid x)\exp\left(\frac{1}{\beta}r(x, y)\right)}_{=\pi(y\mid x)\text{, using Thm.~\ref{thm:main} reparam.}} = 1,
\end{equation}

In other words, $\pi(y\mid x)$ forms a legitimate probability distribution (where probabilities are non-negative and sum to unity). Notably, if we examine Eq.~\ref{eq:op_policy}, we observe that Eq.~\ref{eq:lag_p} corresponds to the partition function for the optimal policy arising from the reward $r(x, y)$. A central idea in the DPO framework is the imposition of specific restrictions on the typically under-constrained Plackett-Luce family of preference models (including Bradley-Terry), maintaining the representational scope of the reward models, while explicitly rendering the optimal policy in Eq.~\ref{eq:op_policy} analytically solvable for any input $x$.

\subsection{Instability of Actor-Critic Algorithms} 

Our framework is equally useful for diagnosing the sources of instability that affect traditional actor-critic approaches in RLHF, like PPO. Adhering to the RLHF methodology, we concentrate on the RL fine-tuning stage detailed in Section \ref{section:prelims}. Here, we establish links with the control as inference approach \cite{levine2018reinforcement} as applied to the constrained RL problem presented in \ref{eq:RL}. Assuming a parameterized policy $\pi_{\theta}(y\mid x)$, our goal is to minimize $\mathbb{D}_{\text{KL}}[\pi_{\theta}(y|x) \mid \mid \pi^*(y\mid x)]$, where $\pi^*$ denotes the optimal policy derived in Eq.~\ref{eq:optimum_model} and shaped by the reward $r_{\phi}(y, x)$. After straightforward algebraic manipulations, this results in the following maximization objective:

\begin{equation}\label{eq:AC}
    \max_{\pi_{\theta}}\mathbb{E}_{\pi_{\theta}(y\mid x)}\bigg[\underbrace{r_{\phi}(x, y) -\beta\log\sum_{y}\pi_\text{ref}(y\mid x)\exp\left(\frac{1}{\beta}r_{\phi}(x, y)\right)}_{f(r_{\phi}, \pi_\text{ref}, \beta)} - \underbrace{\beta\log\frac{\pi_{\theta}(y\mid x)}{\pi_\text{ref}(y\mid x)}}_{\text{KL}}\bigg]
\end{equation}

The objective considered here aligns with that used in previous research 
\citep{ziegler2020finetuning, stiennon2022learning, bai2022training, ouyang2022training}, where the DPO-equivalent reward for the $r_{\phi}$ reward class is applied. In this context, the normalization component within $f(r_{\phi}, \pi_\text{ref}, \beta)$ can be viewed as the soft value function associated with the reference policy $\pi_\text{ref}$. Although this normalization term does not influence the final solution, omitting it can result in high-variance policy gradients, potentially destabilizing training. It is possible to address normalization through the use of a trainable value function, yet such optimization may also pose challenges. Another common approach in the literature involves normalizing rewards with a human completion baseline, which functions as a single-sample Monte-Carlo approximation of the normalization constant. In comparison, the DPO reparameterization generates a reward function that obviates the need for such baselines. 

\section{Experiments}
Here, we empirically assess DPO’s capability for learning policies directly from preference signals. Initially, in a controlled text generation scenario, we investigate the efficiency with which DPO balances maximizing reward and maintaining low KL-divergence with the reference policy, relative to typical preference optimization algorithms such as PPO. Subsequently, we extend our evaluation to larger-scale models and more challenging RLHF benchmarks, including tasks in summarization and dialogue. Our experiments demonstrate that DPO, with minimal hyperparameter adjustment, generally achieves performance on par with or exceeding robust baselines like RLHF with PPO, as well as outperforming strategies that select the best out of $N$ sampled trajectories according to a learned reward. Prior to detailing these findings, we outline our experimental design; further methodological specifics are available in Appendix~\ref{app:exp_details}.

\textbf{Tasks.} Our study investigates three distinct open-ended text generation tasks. In each case, the methods are trained to learn a policy based on a dataset of preferences, represented as $\mathcal{D}=\bigl\{x^{(i)}, y_w^{(i)}, y_l^{(i)}\bigr\}_{i=1}^N$. For the task of \textbf{controlled sentiment generation}, $x$ is defined as the initial segment of a movie review from the IMDb dataset \cite{maas-EtAl:2011:ACL-HLT2011}, and the objective for the policy is to generate an output $y$ exhibiting positive sentiment. To facilitate systematic evaluation, we \textit{construct} preference pairs between candidate continuations using an existing sentiment classifier such that $p(\text{positive}\mid x,y_w)>p(\text{positive}\mid x,y_l)$. The SFT baseline is created by continuing to train GPT-2-large on the positive sentiment reviews within the IMDb training set until the model converges (additional methodological details are found in App~\ref{app:sentiment_details}). In the \textbf{summarization} task, the context $x$ corresponds to a Reddit forum post, with the policy tasked with producing a summary $y$ capturing its central ideas. Adhering to established practices, we utilize the Reddit TL;DR summarization dataset \citep{volske-etal-2017-tl} together with human preference data collected by \citeauthor{stiennon2022learning}. Here, we employ an SFT model adapted on human-curated summaries of forum posts,\footnote{\url{https://huggingface.co/CarperAI/openai_summarize_tldr_sft}} leveraging the TRLX \citep{leandro_von_werra_2023_7790115} framework for RLHF. The corresponding human preference set, annotated by \citeauthor{stiennon2022learning}, was obtained from outputs generated by a different, yet comparably trained, SFT model. In the third scenario, \textbf{single-turn dialogue}, $x$ is a single-turn human prompt, ranging from scientific inquiries to advice requests. The learning objective for the policy is to create a response $y$ that is both useful and engaging; for this, we draw on the Anthropic Helpful and Harmless dialogue dataset \citep{bai2022training}, comprising 170,000 conversation transcripts between human users and automated systems. Each interaction concludes with two responses generated by a large, unidentified language model, accompanied by a label indicating the preferred response according to human judgment. In this domain, since there is no ready-made SFT model, we instead construct one by fine-tuning a generic language model exclusively on the preferred responses.

\textbf{Evaluation.} We employ two distinct evaluation methodologies in our experiments. To rigorously assess how well each algorithm optimizes the reward-maximization objective under constraints, we perform analysis in the controlled sentiment generation context by plotting the reward versus KL-divergence frontier for each algorithm; this analysis is feasible because the underlying reward function (a sentiment classifier) is accessible. In contrast, in real-world scenarios where the ground-truth reward function is unavailable, we turn to a \textit{win rate} metric against a reference policy. Here, GPT-4 serves as an automatic proxy for human evaluators, assessing summary quality in the summarization task and helpfulness of responses in the single-turn dialogue setting. For summarization, baseline comparisons are made to reference summaries from the test set; for dialogue, we use the test set's preferred response as the baseline. Prior research suggests that LMs may outperform existing automatic metrics for evaluation tasks \citep{Chen2023ExploringTU}, but we further validate our use of GPT-4 through a human subject study described in Sec.~\ref{sec:human-judgments}. The results indicate that GPT-4’s ratings are highly aligned with those of human judges, with human agreement rates with GPT-4 that match or surpass inter-annotator agreement.

\textbf{Methods.} Alongside DPO, we benchmark several established strategies for steering language models toward human-preferred outputs. For the summarization task, we apply zero-shot prompting to \textbf{GPT-J} \citep{gpt-j}, while for dialogue we use two-shot prompting with \textbf{Pythia-2.8B} \citep{biderman2023pythia}. Additionally, we assess the \textbf{SFT} approach and a \textbf{Preferred-FT} variant, which undergoes supervised fine-tuning on preferred completions $y_w$—drawn from the SFT model for controlled sentiment and summarization, or from a standard LM for single-turn dialogue. Another approach under consideration is \textbf{Unlikelihood}~\citep{welleck2019neural}, which optimizes the model to maximize the likelihood of $y_w$ and simultaneously \textit{minimize} the likelihood of $y_l$; here, we employ an optional scaling factor $\alpha\in[0,1]$ on the unlikelihood component. We further include \textbf{PPO} \citep{schulman2017proximal}, trained with a reward model learned from preference annotations, as well as \textbf{PPO-GT}, an idealized oracle variant leveraging the ground-truth reward signal in the controlled sentiment setting. For sentiment-related experiments, two PPO-GT implementations are utilized: a standard version \cite{leandro_von_werra_2023_7790115} and a refined variant featuring reward normalization and enhanced hyperparameter optimization for superior results; these improvements are likewise adopted when using PPO with learned reward models. Finally, we introduce the \textbf{Best of $N$} strategy as a baseline: for each input, $N$ candidate outputs are sampled from the SFT model (or from Preferred-FT in dialogue), with the best-performing sample—judged by the learned reward model—returned. Although this approach can achieve strong performance by disentangling the reward model’s quality from the PPO optimization, it becomes computationally expensive for even modest $N$, as it requires multiple samples for each query at evaluation time.

\subsection{How well can DPO optimize the RLHF objective?}

The objective of maximizing reward subject to a KL-divergence constraint, which is central to standard RLHF approaches, serves to balance reward maximization with maintaining similarity to the reference policy. Thus, in assessing different methods, it is essential to jointly consider both the attained reward and the corresponding KL divergence; prioritizing a modest gain in reward at the expense of a significantly greater KL is generally not preferable. Figure~\ref{fig:frontier-tldr-main} presents the trade-off curves between reward and KL for different algorithms in the sentiment classification context. For each method, we perform several training sessions, each employing a distinct value for the policy conservativeness hyperparameter (target KL $\in\{3,6,9,12\}$ for PPO, $\beta \in \{0.05,0.1,1,5\}$, $\alpha\in\{0.05,0.1,0.5,1\}$ for unlikelihood, and varied random seeds for preferred-FT). In total, this results in 22 distinct runs. At intervals of 100 training steps up to convergence, we assess every policy on a suite of held-out prompts, computing both the mean reward from the true reward function and the average KL across entire sequences\footnote{Meaning the aggregate of KL-divergences calculated at each timestep.} with respect to the reference policy $\text{KL}\left(\pi\mid \mid \pi_\text{ref}\right)$. The outcomes indicate that DPO establishes the most effective frontier, yielding superior rewards without incurring high KL. This is a particularly striking finding for a few reasons. Firstly, DPO and PPO are designed to solve the same optimization problem, yet DPO consistently yields better results—its reward/KL tradeoff is always at least as good as PPO's. Secondly, DPO outperforms PPO even when PPO operates with access to ground truth rewards (PPO-GT), achieving a more favorable reward-KL frontier.

\subsection{Can DPO scale to real preference datasets?}
\label{sec:dpo-real-datasets}
We proceed to assess the ability of DPO to fine-tune language models on real-world preference datasets, focusing on tasks such as summarization and single-turn dialogue. For summarization, standard automatic metrics like ROUGE have shown weak alignment with human judgements~\citep{stiennon2022learning}. Earlier studies have demonstrated that fine-tuning language models using PPO with human preference signals yields more effective summary outputs. To benchmark different approaches, we generate model completions for the test set of the TL;DR summarization dataset, then calculate their mean win rate relative to the reference completions in the test data. Completions from each method are sampled across a range of temperatures (0.0 to 1.0), and the win rates are presented in Figure~\ref{fig:frontier-tldr-main} (right). DPO, PPO, and Preferred-FT all start from the same GPT-J SFT model\footnote{\url{https://huggingface.co/CarperAI/openai_summarize_tldr_sft}}. Our experiments show that DPO achieves a win rate near 61\% at a temperature of 0.0, surpassing PPO's best observed rate of approximately 57\% at the same temperature. Furthermore, DPO reaches a higher peak win rate than the best of $N$ method. It is worth noting that we made minimal adjustments to DPO’s $\beta$ hyperparameter, implying that these outcomes may not reflect the upper limit of DPO’s capabilities. Additionally, DPO maintains greater resilience to changes in sampling temperature compared to PPO, whose effectiveness drops to baseline GPT-J levels at higher temperatures. In contrast, Preferred-FT shows negligible improvement over the initial SFT model. Direct comparison of DPO and PPO via human evaluation, discussed in Section~\ref{sec:human-judgments}, further reveals that DPO samples at a temperature of 0.25 were chosen over PPO samples at temperature 0 in 58\% of cases.

For single-turn dialogue tasks, we conduct evaluations using various approaches on the subset of the test partition of the Anthropic HH dataset \citep{bai2022training} that features a single exchange between human and assistant. For the GPT-4 assessments, we utilize the reference completions preferred on the test set to determine the win rates achieved by each method. Since no standard SFT baseline is established for this benchmark, our experiments begin with a pre-trained Pythia-2.8B model. We apply Preferred-FT on the selected completions to build a reference model whose outputs are representative of the in-distribution data, followed by DPO training. Additional comparisons include selecting the best among 128 Preferred-FT completions (noting that the Best of $N$ metric stabilizes at 128 for this scenario, as detailed in Appendix Figure~\ref{fig:best-of-n}) and employing a 2-shot prompt version of the base Pythia-2.8B model. Our findings indicate that DPO matches or outperforms these alternatives when using the optimal temperature setting for each. 

We further assess an RLHF model, which is trained via PPO on the Anthropic HH dataset \footnote{\url{https://huggingface.co/reciprocate/ppo_hh_pythia-6B}} and sourced from a reputable implementation \footnote{\url{https://github.com/CarperAI/trlx/tree/main/examples/hh}}. However, despite attempts with different prompts and sampling temperatures, this RLHF model fails to surpass the performance of the base Pythia-2.8B model. Drawing on our TL;DR experiments and recognizing that both PPO and the Best of 128 strategy optimize for the same reward, we take the Best of 128 as an approximate benchmark for PPO-level effectiveness. Ultimately, DPO emerges as the only efficient method that consistently surpasses the preferred completions in the Anthropic HH dataset, offering results on par with, or superior to, the resource-intensive Best of 128 baseline. Figure~\ref{fig:dialogue-main} further demonstrates that DPO attains peak performance in relatively few training iterations.

\subsection{Generalization to a new input distribution}

To assess the robustness of PPO and DPO to changes in data distribution, we test the policies trained in the Reddit TL;DR summarization setup on an out-of-distribution task: summarizing news articles from the test split of the CNN/DailyMail dataset \citep{nallapati-etal-2016-abstractive}. The same optimal sampling temperatures found for TL;DR (0 and 0.25) are applied. Results are summarized in Table~\ref{tab:ood}. We determined GPT-4’s preference rates between generated summaries and gold references, modifying the prior GPT-4 (C) evaluation prompt used for TL;DR by substituting ``forum post'' with ``news article''. On this dataset shift, DPO policies again show a marked improvement over PPO policies. These findings offer preliminary evidence that DPO generalizes comparably to PPO, despite not leveraging the additional unlabeled Reddit TL;DR prompts that PPO utilizes.

\subsection{Validating GPT-4 judgments with human judgments}
\label{sec:human-judgments}

A human evaluation was carried out to assess the alignment between GPT-4 and human preferences using the outputs of the TL;DR summarization experiment, with two distinct GPT-4 prompting strategies. The \textbf{GPT-4 (S)} (simple) prompt requests the model to identify which summary better conveys key information from the original post. In contrast, the \textbf{GPT-4 (C)} (concise) prompt incorporates a criterion for conciseness, included because the simple prompt can cause GPT-4 to favor longer and sometimes redundant summaries, which diverges from human preferences. Full details of both prompts are provided in Appendix~\ref{app:prompts}. Three system-level comparisons were performed to capture a range of summary quality: the top-performing DPO (temp. 0.25), the lowest-performing PPO (temp. 1.0), and an intermediate SFT (temp. 0.25) policy, all benchmarked against PPO at its best-performing (greedy) temperature. Analysis shows that, for both prompts, GPT-4's preferences align with human raters to a similar degree as inter-human agreement, indicating GPT-4 can serve as a reasonable stand-in for human assessment. Owing to limited human annotation, multiple ratings were only collected for the DPO and PPO-1 pairwise comparisons. In general, the \textbf{GPT-4 (C)} prompt more closely matches human win rates; thus, it is adopted for reporting main results in Section~\ref{sec:dpo-real-datasets}. Further methodological specifics, such as the rater user interface and the roster of human annotators, can be found in Appendix~\ref{app:human-study}.

\section{Discussion}
Preference-based learning constitutes an effective and adaptable approach for developing advanced and well-aligned language models. In this work, we have presented DPO, a straightforward training framework that leverages preferences for language model optimization, circumventing the need for reinforcement learning. Rather than transforming the preference learning objective to fit a traditional RL paradigm and relying on standard RL techniques, DPO establishes a correspondence between language model policies and reward functions, enabling the direct optimization of language models according to human preferences via a simple cross-entropy loss—eliminating the necessity for reinforcement learning or compromising expressiveness. With almost no requirement for hyperparameter tuning, DPO achieves results comparable to or better than current RLHF techniques, including PPO-based strategies. Consequently, DPO substantially lowers the barriers associated with preference-based language model training.

\textbf{Limitations \& Future Work.} Our findings present several avenues for further investigation. For instance, it remains to be seen how policies trained with DPO generalize to out-of-distribution inputs, especially relative to approaches using explicit reward models. Preliminary experiments indicate that DPO-trained policies exhibit generalization properties similar to those achieved by PPO, but a more thorough evaluation is warranted. It would also be valuable to explore whether approaches such as self-labeling with the DPO policy can facilitate effective use of prompts lacking human labels. Another important consideration involves understanding how reward over-optimization arises within the direct preference optimization framework and whether the slight reduction in performance observed in Figure~\ref{fig:dialogue-main}-right provides evidence of this phenomenon. While our empirical analysis has been conducted on models containing up to 6B parameters, extending DPO to contemporary, much larger-scale models presents a compelling research direction. Regarding assessment, we observe that GPT-4-based win rates are sensitive to the phrasing of prompts; thus, future work could focus on methodologies to derive reliable and high-fidelity evaluations from automated systems. Finally, the applicability of DPO extends beyond aligning language models with human preferences, and may be harnessed for training generative models across various modalities.